\section{Discussion}
The full distance $d=5$ magic state cultivation circuit's LER can be calculated exactly on
an everyday device (Apple Macbook M4). The structure exploited is the small syndrome and logical
interface around each non-Clifford double-checking stage (which dominates the non-Clifford cost). 
Pauli propagation reduces each non-Clifford double-checking stage to a response table indexed by the incoming error, 
with stabiliser-equivalent strings merged. Tensor network methods then join these tables to the injection and growth noise, 
whose faults act linearly on the small syndrome-and-logical interface between the stages, so the complete noise average is an exact arithmetic contraction. 
\begin{semiwebaddition}
Pauli propagation also has a graphical interpretation using the Pauli
semiwebs of Ref.~\cite{zxflow}. Fig.~\ref{fig:semiweb-main} shows an
incoming $X_7$ on the $d=3$ double-checking circuit. The red Pauli semiweb in
Fig.~\ref{fig:semiweb-main} shows the Pauli operators on the wires, while the stars mark changes to the $T$-gate phases.
Instead of expanding into signed Pauli terms, this identity keeps a
Pauli string and changes the phases of the circuit it passes through.
Multiplying semiwebs shows where shared Pauli operators and phase changes
cancel. Appendix~\ref{app:semiwebs} defines the notation and gives two
examples.
\end{semiwebaddition}

\begin{figure*}[t]
\centering
\semiwebinput{figures/semiwebs/x7.tikz}
\semiwebcaption{A Pauli semiweb for incoming $X_7$ on the $d=3$
double-checking circuit. Red dashed overlays show $X$ Pauli semiwebs;
unhighlighted wires carry $I$. Stars mark phase changes: $+\pi/2$ at the
entry $T^\dagger_7$, and $-\pi/2$ at the six highlighted exit $T$ gates.
\textbf{Legend (up to normalisation):}
\semiwebCaptionGate{\frac{\pi}{4}} ${}=T$ ($-\pi/4$: $T^\dagger$);
\semiwebCaptionGate{\frac{\pi}{2}} ${}=S$ ($-\pi/2$: $S^\dagger$);
\semiwebCaptionCnot ${}={\rm CNOT}$;
measurement outcome $0$: \semiwebCaptionEffect{Z} ${}=\langle+|$ ($X$ basis),
\semiwebCaptionEffect{X} ${}=\langle0|$ ($Z$ basis);
preparations \semiwebCaptionState{X} ${}=|0\rangle$
and \semiwebCaptionState{Z} ${}=|+\rangle$;
resources \semiwebCaptionPhaseState{\frac{\pi}{4}} ${}=|T\rangle$,
\semiwebCaptionPhaseState{\frac{\pi}{2}} ${}=|Y\rangle$.}
\label{fig:semiweb-main}
\end{figure*}

Our method was loosely inspired by early Pauli-propagation papers based on structured Pauli expansions \cite{lowesa,noisylowesa}. SyQMA~\cite{syqma} previously computed symbolic, noise-averaged LER and inferred circuit fault distances from their leading orders in $p$.
Its Steane-code encoding and verification (akin to magic state cultivation but not exactly) example has $p^2$ scaling, under a different circuit and noise conventions.

In addition, the fault-count coefficients $a_k,e_k$ are related to the per-stratum acceptance and accepted-error probabilities that Clifft's stratified importance sampling
estimates~\cite{clifft} and that the Pauli-frame reweighting of~\cite{pfsr} organises by fault number. Here every stratum is summed
exactly rather than sampled. In~\cite{pfsr}, they restricts their $d=5$ sampling
to $k\ge5$ faults on the premise that fewer than $d$ faults cannot yield an
accepted logical error. For the circuits studied here that premise fails,
since $e_3$ and $e_4$ are non-zero and supply about a third of the accepted errors at $p=10^{-3}$ when faults are counted over all 3,564 physical locations, and four fifths when only the 1,996 locations with a non-identity response are counted, Appendix~\ref{app:circuits}.

\begin{figure}[h]\centering
\includegraphics[width=.98\linewidth]{figures/d5_t_vs_s.pdf}
\caption{$d=5$ cultivation with unitary injection, without escape.
The {\color{red}$T$ curve} uses exact series coefficients through $x^{10}$. The {\color{blue}$S$-proxy data (blue)} and {\color{green}green curve} (that retains detector-error-model terms of at most weight five) are both from~\cite{msc}.}
\label{fig:d5-s-proxy}
\end{figure}

Fig.~\ref{fig:d5-s-proxy} compares the $d=5$ $T$-state series with the
unitary-injection $S$-gate proxy of~\cite{msc}.
At $p=10^{-3}$, the proxy is $\sim 9.6$ times lower (nearly one order of
magnitude). 
Using the proxy to infer a quantum algorithmic estimate will produce an over-optimistic resource count, as the proxy may substantially underestimate the LER of actual $T$-state cultivation.

The leading orders of the series expansion reveal the LER scales as $p^{2}$ at $d=3$ and $p^{3}$ at $d=5$ (see Acknowledgements), not as the intended $p^{3}$ and $p^{5}$ the code distances suggest. 
This explains the apparent distance-degrading behaviour observed in the companion code of~\cite{wan}, albeit with a slightly different injection scheme. 
In a Clifford circuit, detectors are equivalent to closed Pauli webs of its ZX diagram. In the $S$-gate proxy this holds for all 20 declared $d=3$ detectors, but in the $T$-gate circuit only 16 of them are Pauli webs (93 of 107 at $d=5$), and a fault pattern that the $T$-gate circuit accepts can be rejected by the proxy only through the other detectors (Appendix~\ref{app:detector-webs}).
Appendix~\ref{app:other-cultivation} provides exact series coefficients for other flavours of magic state cultivation such as: $\mathbb{RP}^{2}$~\cite{RP2_prxq}, fold-transversal~\cite{fold_msc_prxq} and repaired cultivation~\cite{chan2026diagnosingrestoringdegradedfault}, with comparisons against available numerical results.
Extensions to the escape and grafting stages of cultivation~\cite{msc} are left for future work.
